\documentclass[11pt]{article}
\usepackage[a4paper,margin=1in]{geometry}
\usepackage{amsmath,amssymb,amsthm,mathtools}
\usepackage{graphicx}
\usepackage[hidelinks]{hyperref}
\usepackage{microtype}

\newtheorem{theorem}{Theorem}
\newtheorem{lemma}[theorem]{Lemma}
\theoremstyle{remark}
\newtheorem{remark}[theorem]{Remark}

\newcommand{\C}{\mathbb C}
\newcommand{\M}{\mathbb M}
\newcommand{\MR}{\mathcal{MR}}
\newcommand{\J}{\mathcal J}
\newcommand{\rank}{\operatorname{rank}}

\title{A sharp rank-eight extremal marginal state in dimension six}
\author{Candidate partial solution packet for arXiv:2412.10041}
\date{12 August 2026}

\begin{document}
\maketitle

\begin{center}
\fbox{\parbox{0.91\textwidth}{\textbf{Status.}
Candidate substantial partial solution, computer-assisted but exact.
We settle the first equal-dimension case omitted by the source:
\(\MR(I_6/6,I_6/6)=8\).  The general construction problem remains open.}}
\end{center}

\section{Source question}

For states \(\rho_1\in\M_{d_1}\) and \(\rho_2\in\M_{d_2}\), let
\(\mathcal C(\rho_1,\rho_2)\) be the convex set of bipartite states with
these two marginals.  Parthasarathy's bound says that every extreme
\(\rho\in\mathcal C(\rho_1,\rho_2)\) satisfies
\[
 \rank\rho\le
 \left\lfloor\sqrt{d_1^2+d_2^2-1}\right\rfloor .
\]
Devendra, Dey, and Dey ask for constructions attaining this bound.  Their
paper handles \(d_1=d_2=3,4,5,9,12\) and certain multiples of five, but not
\(d_1=d_2=6\).  In that case the bound is
\[
 \left\lfloor\sqrt{2\cdot6^2-1}\right\rfloor
 =\lfloor\sqrt{71}\rfloor=8.
\]

\begin{figure}[ht]
 \centering
 \includegraphics[width=0.93\textwidth]{figures/source_open_problem_page.png}
 \caption{The source's concluding construction problem, on PDF page 17.}
\end{figure}

\section{Proof intuition}

Under the Choi--Jamio\l kowski correspondence, it is enough to construct a
bistochastic completely positive map on \(\M_6\) with eight Kraus operators.
The marginal constraints ask for
\[
 \sum_j U_j^*U_j=I_6=\sum_j U_jU_j^*.
\]
Extremality asks that the \(64\) ordered pairs
\[
 \bigl(U_i^*U_j,\;U_jU_i^*\bigr),\qquad 1\le i,j\le8,
\]
be linearly independent in \(\M_6\oplus\M_6\).

The construction makes every \(U_j\) monomial.  Across every input column,
and also across every output row, the eight nonzero moduli are exactly
\(1,2,\ldots,8\).  Both marginal identities then follow from the same
one-line sum of squares.  The remaining independence question is finite
linear algebra over the Gaussian integers; one nonzero determinant modulo a
prime is an exact certificate.

\section{The construction}

Let \(e_1,\ldots,e_6\) be the standard basis of \(\C^6\).  Define eight
matrices by listing their columns:
\small
\begin{align*}
W_1&=[-8e_3,-8ie_4,-8ie_5,6ie_2,e_1,4e_6],\\
W_2&=[7e_1,-7e_5,-3e_2,-5e_6,-6e_4,-7ie_3],\\
W_3&=[4e_3,2e_1,2e_6,-7e_2,7ie_4,2ie_5],\\
W_4&=[-3e_4,3e_5,5ie_3,4e_1,-8ie_6,8ie_2],\\
W_5&=[2e_4,6ie_6,4ie_2,-8ie_1,4ie_5,ie_3],\\
W_6&=[6e_1,5e_4,6e_3,-e_6,5ie_5,5ie_2],\\
W_7&=[5ie_1,-e_2,-ie_4,2ie_3,3ie_6,-6e_5],\\
W_8&=[-e_5,-4ie_4,7e_6,-3e_1,-2e_2,3ie_3].
\end{align*}
\normalsize

\begin{lemma}\label{lem:certificate}
The matrices above satisfy
\[
 \sum_{j=1}^8W_j^*W_j=204I_6
 \quad\text{and}\quad
 \sum_{j=1}^8W_jW_j^*=204I_6.
\]
Moreover, the \(64\) pairs
\[
 \left\{\bigl(W_i^*W_j,W_jW_i^*\bigr):1\le i,j\le8\right\}
\]
are linearly independent over \(\C\).
\end{lemma}

\begin{proof}
Every \(W_j\) is monomial.  Inspection of the displayed columns shows that,
for each fixed input coordinate, the eight nonzero moduli form the multiset
\(\{1,\ldots,8\}\).  After transporting each entry to its output row, the
same is true for every fixed output coordinate.  Off-diagonal entries in
both sums vanish, while each diagonal entry equals
\[
 1^2+2^2+\cdots+8^2=204.
\]

For independence, form a \(72\times64\) Gaussian-integer matrix \(A\).
Its columns are indexed lexicographically by \((i,j)\).  The first \(36\)
rows are the entries of \(W_i^*W_j\), in row-major order, and the final
\(36\) rows are the entries of \(W_jW_i^*\), again in row-major order.
Let \(D\) be the leading \(64\times64\) minor of \(A\); thus \(D\) uses all
first-product rows and the first \(28\) second-product rows.

Reduce \(\mathbb Z[i]\) to \(\mathbb F_{101}\) by sending \(i\) to \(10\).
This is a ring homomorphism because \(10^2\equiv-1\pmod{101}\).  Exact
Gaussian elimination gives
\[
 \det D\equiv19\pmod{101}.
\]
Hence \(\det D\ne0\) in \(\mathbb Z[i]\), so \(A\) has column rank \(64\)
over \(\C\).  The claimed pairs are therefore linearly independent.  The
included verifier reconstructs all matrices and performs both normalization
checks and this modular determinant computation using integer arithmetic
only.
\end{proof}

\begin{theorem}\label{thm:main}
There is an extreme point
\[
 \rho\in\mathcal C(I_6/6,I_6/6)
\]
of rank \(8\).  Consequently,
\[
 \boxed{\MR(I_6/6,I_6/6)=8.}
\]
\end{theorem}

\begin{proof}
Put \(U_j=W_j/\sqrt{204}\) and define
\[
 \Phi(X)=\sum_{j=1}^8U_j^*XU_j,\qquad X\in\M_6.
\]
The two identities in Lemma~\ref{lem:certificate} say that \(\Phi\) is
unital and trace preserving.  The extremality criterion recalled as
Theorem 2.12 in the source says that \(\Phi\) is extreme among such maps
exactly when the pairs
\((U_i^*U_j,U_jU_i^*)\) are linearly independent.  Lemma
\ref{lem:certificate} gives this independence.

It also forces the \(U_j\)'s themselves to be linearly independent.  Indeed,
if \(\sum_jc_jU_j=0\), fixing a nonzero \(U_k\) yields simultaneous
relations
\[
 \sum_jc_jU_k^*U_j=0,\qquad
 \sum_jc_jU_jU_k^*=0,
\]
and pairwise independence gives every \(c_j=0\).  Thus the Kraus
decomposition is minimal and the Choi rank of \(\Phi\) is \(8\).

Now set \(\rho=\J(\Phi)/6\).  The Choi correspondence gives
\(\rho\in\mathcal C(I_6/6,I_6/6)\), preserves extremality, and identifies
\(\rank\rho\) with the Choi rank.  Hence \(\rho\) is extreme of rank \(8\).
Parthasarathy's upper bound is also \(8\), proving equality.
\end{proof}

\section{Exact verification}

Run
\begin{verbatim}
python verify_exact.py
\end{verbatim}
from the packet directory.  The script uses no numerical linear algebra or
external package.  It verifies the two identities with \(204I_6\), builds
the \(64\times64\) minor, and returns determinant \(19\) in
\(\mathbb F_{101}\).  The certificate is deliberately stronger than a
floating-point rank computation.

\section{Upgrade attempts, limitations, and novelty}

The first construction attempt used arbitrary Kraus matrices and operator
scaling, but it would have left irrational numerical normalization without
a compact certificate.  The successful upgrade replaced it by monomial
Gaussian-integer matrices whose moduli form two simultaneous edge colourings;
this makes both marginals exact before the rank test.

A second upgrade generalized the search to other omitted dimensions.  Exact
full-rank modular candidates appeared immediately in several cases, strongly
suggesting a generic phenomenon.  They are not claimed here because a
uniform proof that the relevant determinant polynomial is nonzero for every
dimension is still missing.  A dimension-count/transversality argument alone
is insufficient: the \(d=2\) bistochastic case is a genuine exception despite
the naive count.  Likewise, tensor-product propagation gives only finite
ranges because Choi ranks multiply while the sharp integer bound follows
\(\lfloor\sqrt2\,d\rfloor\).

Thus Theorem~\ref{thm:main} completely settles the first omitted dimension
but not the source's general algorithmic question.

A bounded primary-source search through 12 August 2026 used the exact source
title and arXiv id, the phrase extremal marginal state dimension 6 rank 8,
the phrase bistochastic channel Choi rank 8, and the Parthasarathy sharpness
question.  It found the source's cases and older general extremality
criteria, but no construction for \(d=6\) and no later general resolution.
Novelty confidence is moderate; specialist review is recommended before any
priority claim.

\begin{thebibliography}{9}
\bibitem{DDD}
R.~Devendra, P.~Dey, and S.~Dey,
\emph{On the rank of extremal marginal states},
arXiv:2412.10041 (2024).
\end{thebibliography}

\end{document}
